\chapter{Smoothing, coverage, and the connecting sequence}\label{u:sec:connections}

The goal of this section is to prove \cref{u:prop:witness}, which excludes
balanced cuts with small stationary flow on the event
$\|W(g)\|_{\mathrm{op}}\le12$. We use the fine entropy estimate of
\cref{u:prop:fine-kl} to obtain coverage of both sides of a cut, and the
coarse domination estimate of \cref{u:prop:coarse} to construct a sequence
that avoids its enlarged boundary. These are the smoothing and connecting
sequence arguments of \cite[Sections~5--6, Propositions~6--7]{oai-typical-upper},
with the estimates at the fine and coarse horizons established in
\cref{u:sec:fine,u:sec:coarse}. In \cref{u:sec:proof}, the witness estimate
and \cref{u:lem:goe-norm} will imply \cref{thm:upper}.

Throughout, we fix $\beta>1$ and the exponent parameters before taking
$n\to\infty$, as in \cref{u:sec:parameters}. We retain the primitive disorder,
completed Hamiltonian, and joint law $\Pi_n$ of \cref{u:sec:primitive},
and the source and replay constructions of \cref{u:sec:proposal}.

\section{Smoothing and coverage}\label{u:sec:smoothing}

We first record the transport inequality for disorder conditioned on a
spin. Under $\Pi_n$, the spin marginal is $\Unif_n$: this follows from the
gauge transformation $g_{ij}\mapsto x_ix_jg_{ij}$ and summation over spins.
Conditional on $x$, primitive disorder has density
\begin{equation}\label{u:eq:conditional-density}
 p_x(g)=2^n(2\pi)^{-n^2/2}
       \exp\{-\|g\|_F^2/2+\beta H^c_g(x)-F^c(g)\}.
\end{equation}
Since $F^c$ is a log-sum-exp of linear functions and $H^c_g(x)$ is
linear in $g$, we have that the negative logarithm of $p_x$ is
$1$-strongly convex.
We write $W_2$ for the quadratic Wasserstein distance for the Euclidean
norm on the primitive matrix space, that is, the Frobenius norm.

\begin{lemma}\label{u:lem:transport}
For every $x\in\Sigma_n$ and every probability measure $\eta$ on the
primitive matrix space,
\begin{equation}\label{u:eq:conditional-transport}
 W_2(\eta,p_x)^2\le2\KL(\eta\Vert p_x).
\end{equation}
\end{lemma}

\begin{proof}
This is the transport inequality for measures with a strongly convex potential; see
\cite[Corollary~3.10]{bogachev-kolesnikov-medvedev-2005} and its use in
\cite[Section~5]{oai-typical-upper}. We recall the finite-dimensional
argument. For a smooth positive density $\eta$, let $T$ be the optimal
transport map from $p_x$ to $\eta$, which is the gradient of a convex function, and let $V_x$ be the
negative logarithm of $p_x$. By the change-of-variables formula, we have
\[
 \KL(\eta\Vert p_x)
 =\int\bigl(V_x(T(g))-V_x(g)-\log\det DT(g)\bigr)p_x(g)\dd g.
\]
By strong convexity and $-\log\det DT\ge\tr(I-DT)$, we have
\[
 \KL(\eta\Vert p_x)
 \ge\frac12\int\|T(g)-g\|_F^2p_x(g)\dd g.
\]
In the last step, we used integration by parts to cancel the linear terms.
For a general measure with finite entropy, we obtain the inequality by smooth
truncation and regularization, followed by lower semicontinuity. If the
entropy is infinite, there is nothing to prove.
\end{proof}

For every $L$-Lipschitz linear function $f$ and every $t\in\R$, the variational formula for relative entropy and \cref{u:lem:transport} imply that
\begin{equation}\label{u:eq:conditional-linear-mgf}
 \E_{p_x}\exp\{t(f-\E_{p_x}f)\}\le e^{t^2L^2/2}.
\end{equation}
Indeed, $\E_\eta f-\E_{p_x}f\le L W_2(\eta,p_x)$, and maximizing
$|t|L\sqrt{2K}-K$ over $K\ge0$ proves
\eqref{u:eq:conditional-linear-mgf}.

Let $\delta=n^{-\ex{D}}$, and write the fine output as
$Q_F(\dd g,x)=\gamma_n(\dd g)q_F(g,x)$. We take the version of $q_F$
given by the replay integral, with the conventions of
\cref{u:sec:primitive}, and define
\begin{equation}\label{u:eq:smoothed-kernel}
 \nu^{\rm sm}_g(x)=\int
 q_F(\sqrt{1-\delta}\,g+\sqrt\delta\,z,x)\,\gamma_n(\dd z).
\end{equation}
For every $g$, we have that $\nu^{\rm sm}_g$ is a probability measure
on the spins.

\begin{lemma}\label{u:lem:smoothing}
Let $g_\star,g_\circ$ be independent with law $\gamma_n$, and set
$g_{\rm base}=\sqrt{1-\delta}\,g_\star+\sqrt\delta\,g_\circ$. Then
\begin{equation}\label{u:eq:reverse-entropy}
 \E_{g_\star,g_\circ}\KL(\pi_{g_\star}\Vert\nu^{\rm sm}_{g_{\rm base}})
 \le\frac{\beta^2(n-1)}2\delta+\delta^{-1}\KL(Q_F\Vert\Pi_n).
\end{equation}
In particular, if $\ex{u}_{\rm E}<\min(\ex{D},\ex{k}_{\rm f}-\ex{D})$,
the outer event
\begin{equation}\label{u:eq:coverage-event}
 \mathcal C_n=\bigl\{(g_\star,g_\circ):
 \KL(\pi_{g_\star}\Vert\nu^{\rm sm}_{g_{\rm base}})
                         \le n^{1-\ex{u}_{\rm E}}\bigr\}
\end{equation}
has complementary probability at most
$C_\beta(n^{\ex{u}_{\rm E}-\ex{D}}+
 n^{\ex{u}_{\rm E}-\ex{k}_{\rm f}+\ex{D}})=o(1)$.
\end{lemma}

\begin{proof}
We first compare the joint laws after smoothing, conditioning on the spin.
We then compare the Gibbs laws at the target and its smoothing base.
This is the reverse comparison of \cite[Section~5, Proposition~6]{oai-typical-upper};
we use Gaussian integration by parts for the second comparison.

We disintegrate $\Pi_n$ and $Q_F$ over the spin, whose marginal is uniform
under both laws. For each spin, by \eqref{u:eq:conditional-transport},
we can couple the conditional laws of the disorder with squared cost at
most twice their forward relative entropy. Under the channel
$g\mapsto\sqrt{1-\delta}\,g+\sqrt\delta\,Z$, the relative entropy between
the Gaussian output laws at two inputs is $(1-\delta)\|g-g'\|_F^2/(2\delta)$ in either
direction. By convexity of relative entropy, applied to the coupled mixture and
averaged over the common spin, we have
\begin{equation}\label{u:eq:smoothed-joint-entropy}
 \KL(\Pi_n^{\rm sm}\Vert Q_F^{\rm sm})
                 \le\delta^{-1}\KL(Q_F\Vert\Pi_n).
\end{equation}
This comparison for Gaussian convolutions also appears in
\cite[Lemma~11]{lugosi-neu-2022}. Both smoothed measures have disorder
marginal $\gamma_n$. By Gaussian reversibility, we identify the
conditional spin law of $Q_F^{\rm sm}$ with \eqref{u:eq:smoothed-kernel},
and we denote that of $\Pi_n^{\rm sm}$ by $\bar\pi_g$. Under the joint law of $(g_\star,g_{\rm base})$, expanding the logarithms
of the probabilities on the finite spin space gives
\begin{equation}\label{u:eq:entropy-decomposition}
 \begin{split}
 \E\KL(\pi_{g_\star}\Vert\nu^{\rm sm}_{g_{\rm base}})
 &=\E\KL(\pi_{g_\star}\Vert\bar\pi_{g_{\rm base}})\\
 &\quad+\E\KL(\bar\pi_{g_{\rm base}}\Vert\nu^{\rm sm}_{g_{\rm base}}).
 \end{split}
\end{equation}
Using the same decomposition with any alternative kernel at $g_{\rm base}$,
we find that the first term is at most
$\E\KL(\pi_{g_\star}\Vert\pi_{g_{\rm base}})$. When intermediate expressions may be infinite, we justify these
identities by finite approximations, or by the variational formula for relative entropy and monotone
approximation. By \cref{u:prop:fine-kl} and the following calculation,
the upper bounds are finite.

Since $g_\star$ and $g_{\rm base}$ have the same marginal law, the expected
difference of log partition functions vanishes. Using the independence of
$g_\circ$, we obtain
\begin{equation}\label{u:eq:gibbs-smoothing-cost}
 \begin{split}
 \E\KL(\pi_{g_\star}\Vert\pi_{g_{\rm base}})
 &=\beta(1-\sqrt{1-\delta})
   \E_{g_\star}\sum_x\pi_{g_\star}(x)H^c_{g_\star}(x)\\
 &\le\tfrac12\beta^2(n-1)(1-\sqrt{1-\delta})
 \le\tfrac12\beta^2(n-1)\delta.
 \end{split}
\end{equation}
For the middle inequality, Gaussian integration by parts gives the
off-diagonal contribution
$(\beta/n)\sum_{i<j}(1-\langle x_ix_j\rangle_{g_\star}^2)$;
the spin-independent diagonal has mean zero. Since the derivatives that occur here are bounded spin moments and
Gaussian first moments are finite, the integrations are justified. Combining \eqref{u:eq:smoothed-joint-entropy}--\eqref{u:eq:gibbs-smoothing-cost}
proves \eqref{u:eq:reverse-entropy}. The final assertion follows from
\cref{u:prop:fine-kl} and Markov's inequality.
\end{proof}

\begin{lemma}\label{u:lem:coverage}
On $\mathcal C_n$, every spin has positive
$\nu^{\rm sm}_{g_{\rm base}}$ probability, and, simultaneously for every
$S\subseteq\Sigma_n$,
\begin{equation}\label{u:eq:all-set-coverage}
 \pi_{g_\star}(S)\ge\epsilon_n=n^{-\ex{d}}
 \quad\Longrightarrow\quad
 \nu^{\rm sm}_{g_{\rm base}}(S)
 \ge\exp\{-n^{\ex{d}}(n^{1-\ex{u}_{\rm E}}+\log2)\}.
\end{equation}
If $\ex{e}+\ex{d}<\ex{u}_{\rm E}$, two independent draws from
$\nu^{\rm sm}_{g_{\rm base}}$ cross every $\epsilon_n$-balanced cut with
probability at least
\begin{equation}\label{u:eq:crossing-probability}
 \varpi_n=\exp\{-2n^{\ex{d}}(n^{1-\ex{u}_{\rm E}}+\log2)\}
                       \ge e^{-2n^{1-\ex{e}}}
\end{equation}
eventually.
\end{lemma}

\begin{proof}
Finite relative entropy and the strict positivity of $\pi_{g_\star}(x)$
imply that $\nu^{\rm sm}_{g_{\rm base}}(x)>0$ for every spin. We pass to the partition $\{S,S^c\}$, which does not increase relative
entropy. If $a=\pi_{g_\star}(S)$ and
$b=\nu^{\rm sm}_{g_{\rm base}}(S)$, the binary relative entropy is at least
$a\log(1/b)-\log2$, since binary entropy is at most $\log2$.
We have proved \eqref{u:eq:all-set-coverage}. Applying it to both sides
of a balanced cut and multiplying the independent probabilities gives
\eqref{u:eq:crossing-probability}; the strict power gap also absorbs the
$\log2$ term. The implication holds for every set once the outer pair is
fixed, including a cut selected from this pair.
\end{proof}

\section{The witness event and the enlarged boundary}\label{u:sec:witness}

We now select a balanced cut with small flow, whenever one exists, and
bound the Gibbs mass of its enlarged boundary. The selection follows
\cite[Section~6, Proposition~7]{oai-typical-upper}. The estimate below uses
the conditional disorder law to control sums of local fields over small
sets of sites.

For $\epsilon>0$ and $\phi>0$, we define the witness event
\begin{equation}\label{u:eq:witness-event}
 \begin{split}
 \mathcal W_n(\epsilon,\phi)=\{g:\ &\|W(g)\|_{\mathrm{op}}\le12,\\
 &\exists S\subseteq\Sigma_n:\epsilon\le\pi_g(S)\le1-\epsilon,
                         \ \mathcal Q_g(S,S^c)\le\phi\}.
 \end{split}
\end{equation}
Fix an ordering of the subsets of $\Sigma_n$, and let $S_g$ be the first
cut satisfying these conditions. The cut $S_g$ depends only on $g$, and any fixed order
works for the argument. The mass and flow of each cut are continuous
finite sums in $g$, so this selection is measurable, as in the conventions
of \cref{u:sec:primitive}. The selected sets defined below are empty off
the witness event.

Let
\[
 \epsilon_n=n^{-\ex{d}},\qquad \phi_n=e^{-n^{1-\ex{q}}},\qquad
 R_n=\lceil n^{1-\ex{r}}\rceil,\qquad
 A_n^{\rm fld}=n^{1-\ex{u}_{\rm F}}.
\]
We write $d_H(x,y)=\#\{i\in[n]:x_i\ne y_i\}=\|x-y\|_2^2/4$
for the Hamming distance between $x,y\in\Sigma_n$, and
$d_H(x,A)=\min_{y\in A}d_H(x,y)$, with
$d_H(x,\varnothing)=+\infty$.
We write $\partial S$ for the vertex boundary that contains both endpoints
of every edge across the cut, and define
$\partial_R S=\{x:d_H(x,\partial S)\le R\}$. Let $\mathcal V_R$ be the family of differences of the indicator vectors
of two subsets of $[n]$ with at most $R$ elements each. Then $v_i\in\{-1,0,1\}$ and
$\|v\|_1=\|v\|_2^2\le2R$. The family contains both signs of every
indicator of at most $R$ sites, and
\begin{equation}\label{u:eq:field-direction-count}
 |\mathcal V_R|\le[(R+1)(n+1)^R]^2.
\end{equation}
Define
\[
 \mathrm{Bad}_{R,A}(g)=\Bigl\{x:\exists v\in\mathcal V_R,
                         \ \sum_i v_i(W(g)x)_i\ge A\Bigr\}.
\]
The excluded set is
\begin{equation}\label{u:eq:excluded-set}
 \mathcal E_n(g)=\partial_{R_n}S_g
           \ \cup\ \mathrm{Bad}_{R_n,A_n^{\rm fld}}(g)
 \quad\text{on }\mathcal W_n(\epsilon_n,\phi_n),
\end{equation}
and is empty off that event. It contains the enlarged boundary of the selected cut and the spins with large sums of local fields.

\begin{lemma}\label{u:lem:boundary}
Suppose that
\begin{equation}\label{u:eq:boundary-conditions}
 0<\ex{q}<\ex{u}_{\rm F}<\ex{r}<1,\qquad
 2\ex{u}_{\rm F}-\ex{r}<\ex{q},\qquad \ex{q}<1/2.
\end{equation}
Then, eventually,
\begin{equation}\label{u:eq:annealed-boundary-mass}
 \E_{\gamma_n}\pi_g(\mathcal E_n(g))
                         \le e^{-\frac12 n^{1-\ex{q}}}.
\end{equation}
\end{lemma}

\begin{proof}
We first bound the annealed mass of $\mathrm{Bad}_{R,A}$ by conditioning
on the spin. We then compare Gibbs weights along a short path from an
enlarged boundary vertex to the boundary.

Fix a spin $x$. Recall that, under the joint law $\Pi_n$, the conditional
density of the disorder given this spin is
\eqref{u:eq:conditional-density}. For fixed $v$, the linear functional
$f_v(g)=\sum_i v_i(W(g)x)_i$ has coefficients
$a_{ij}=(v_ix_j+v_jx_i)/\sqrt{2n}$, so
\[
 \|a\|_2^2=\|v\|_2^2+(\textstyle\sum_i v_ix_i)^2/n
                         \le2\|v\|_2^2\le4R.
\]
By integration by parts under \eqref{u:eq:conditional-density}, we have
$|\E_{p_x}f_v|\le2\beta\|v\|_1\le4\beta R$. Indeed, the derivative of
$\beta H^c_g(x)-F^c(g)$ in the $(i,j)$ entry is
$\beta(x_ix_j-\langle y_iy_j\rangle_g)/\sqrt{2n}$, whose absolute value
is at most $2\beta/\sqrt{2n}$. The density has Gaussian tails up to a linear exponential factor; by
truncation, this justifies the integration by parts, including its
boundary terms. Using
\eqref{u:eq:conditional-linear-mgf} and a union bound over
$\mathcal V_R$, we find that, for $A\ge4\beta R$,
\begin{equation}\label{u:eq:sparse-field-tail}
 \E_{\gamma_n}\pi_g(\mathrm{Bad}_{R,A}(g))
 \le[(R+1)(n+1)^R]^2
                    \exp\{-(A-4\beta R)^2/(8R)\}.
\end{equation}
When $R=R_n$ and $A=A_n^{\rm fld}$, the inequality
$\ex{u}_{\rm F}<\ex{r}$ implies that $A\ge8\beta R$ eventually. We also
have that the logarithm of \eqref{u:eq:field-direction-count} is $o(A^2/R)$.
Then \eqref{u:eq:sparse-field-tail} implies the upper bound
$\exp(-c_\beta n^{1+\ex{r}-2\ex{u}_{\rm F}})$ eventually.
By \eqref{u:eq:boundary-conditions}, the power of $n$ in this exponent is
strictly larger than $1-\ex{q}$.

Next, fix a witness matrix and its selected cut. On
$\|W(g)\|_{\mathrm{op}}\le12$, we find by polarization that the change
in energy when one spin is flipped is at most $24\sqrt n$. Then every
flip rate is at least
$(1/2)e^{-24\beta\sqrt n}$. By reversibility, the flow in each direction
across the selected cut is at most $\phi_n$. We conclude that its two-sided
boundary has Gibbs mass at most $4\phi_n e^{24\beta\sqrt n}$.
If $x\notin\mathrm{Bad}_{R,A}(g)$ and $d_H(x,y)\le R$, expansion about
$x$ gives
\begin{equation}\label{u:eq:sparse-energy-change}
 |H^c_g(y)-H^c_g(x)|\le2A+24R.
\end{equation}
To see this, let $I$ be the set of flipped sites. Both $-x\mathbf 1_I$
and its negative belong to $\mathcal V_R$, since $|I|\le R$. Their linear
fields are bounded by $A$ in both signs, and the quadratic correction is
at most $2\|W(g)\|_{\mathrm{op}}R$.
For each vertex of the enlarged boundary, we select a nearest boundary
vertex, using a fixed order to break ties. By the bound on the size of a
Hamming ball, each boundary vertex is selected at most
$(R+1)(n+1)^R$ times. Comparing Gibbs weights by
\eqref{u:eq:sparse-energy-change}, we obtain
\begin{equation}\label{u:eq:boundary-weight-comparison}
 \begin{split}
 \pi_g(\mathcal E_n(g))
 &\le4(R_n+1)(n+1)^{R_n}\phi_n
       e^{\beta(2A_n^{\rm fld}+24R_n+24\sqrt n)}\\
 &\quad+\pi_g(\mathrm{Bad}_{R_n,A_n^{\rm fld}}(g)).
 \end{split}
\end{equation}
The inequality also holds off the witness event, since its left side
vanishes there. In the first term, the logarithms of the positive factors other than
$\phi_n$ have powers $1-\ex{r}$ (with a factor $\log n$), $1-\ex{u}_{\rm F}$,
and $1/2$, all smaller than $1-\ex{q}$. Averaging
\eqref{u:eq:boundary-weight-comparison} and using
\eqref{u:eq:sparse-field-tail}, we obtain
\eqref{u:eq:annealed-boundary-mass} after absorbing the sum and constants.
\end{proof}

We impose a pressure cutoff only at the outer target. We define
\begin{equation}\label{u:eq:target-pressure-event}
 \begin{split}
 \mathcal H_n&=\{g:F^c(g)\le nP_\beta+n^{1-\ex{z}_{\rm P}}\},\\
 \overline{\mathcal E}_n(g)&=
 \begin{cases}
 \mathcal E_n(g),&g\in\mathcal H_n,\\
 \varnothing,&g\notin\mathcal H_n.
 \end{cases}
 \end{split}
\end{equation}
The gradient of $F^c$ has norm at most $\beta\sqrt{n/2}$.
By Gaussian concentration and $\E F^c\le nP_\beta$
(\cref{inp:ob-concentration,inp:parisi,lem:model-completion}), we have
\begin{equation}\label{u:eq:target-pressure-tail}
 \gamma_n(\mathcal H_n^c)
 \le e^{-c_\beta n^{1-2\ex{z}_{\rm P}}}=o(1)
                         \qquad(0<\ex{z}_{\rm P}<1/2).
\end{equation}
This is also \cref{u:lem:pressure-tail}. Each set is measurable by the
conventions of \cref{u:sec:primitive}; the disorder used in each coarse
proposal is not restricted by $\mathcal H_n$.

\section{The connecting sequence}\label{u:sec:sequence}

We construct the sequence on one product space, following
\cite[Section~6, Proposition~7]{oai-typical-upper}. At the two endpoints
we continue the coarse proposal to the fine horizon. At every other
vertex we use only the coarse proposal.

For the common fine schedule, let $\mu^{\rm aux}_n$ be the independent
product law of all $D_i$ Wiener paths to $V_{\ex{f}}$, auxiliary $B_i^0$
paths, used bridge paths through $N_{\ex{f}}$, and full completion vectors
$C_j\sim N(0,I_n)$ through $N_{\ex{f}}$. We may retain the unused bridge coordinates independently; this gives
the same finite experiment. At both horizons, the matrix has dimension $n$ and the initial frame is
the identity. By
\cref{u:lem:assembly}, these are the replay inputs.

Under a single product probability $\PP_n$, we sample independently
\[
 g_\star,g_\circ,g_1,g_2\sim\gamma_n,\qquad
 a_1,a_2\sim\mu^{\rm aux}_n,
\]
and independent arrays of uniform thresholds at each grid point, together
with two arrays for the fine endpoints. For $0\le\theta\le\pi/2$, we define
\begin{equation}\label{u:eq:coupled-rotation}
 \begin{split}
 g_{\rm base}&=\sqrt{1-\delta}\,g_\star+\sqrt\delta\,g_\circ,\\
 g_\theta&=\sqrt{1-\delta}\,g_{\rm base}
                +\sqrt\delta(\cos\theta\,g_1+\sin\theta\,g_2),\\
 a_\theta&=\cos\theta\,a_1+\sin\theta\,a_2.
 \end{split}
\end{equation}
The rotation of the auxiliary inputs is defined by blockwise addition of
vectors and paths, which preserves the zero endpoints of bridges. The uniform thresholds are not
rotated. We use the coarse replay at every angle and continue it to the fine
horizon at the two endpoints. Let
\[
 \ell_n=\lceil(\pi/2)e^{n^{\ex{g}_{\rm rot}}}\rceil,
 \qquad \theta_j=\frac{\pi j}{2\ell_n}\quad(0\le j\le\ell_n).
\]
We denote the coarse spins by $X_j^B$, their coarse magnetizations by
$M_j$, and the fine endpoint spins by $X_0^F,X_1^F$.
For a grid angle $\theta=\theta_j$, we write $X_\theta^B=X_j^B$,
and let $U_j^B$ be its coarse threshold array. We define
\[
 E_{B,\theta_j}
   =\{(g_{\theta_j},a_{\theta_j}|_B,U_j^B)\in E_B\},
\]
where $E_B$ is the event of \cref{u:prop:coarse} and
$a|_B$ restricts the $D$ paths to $V_{\ex{b}}$ and retains only
the auxiliary cell, bridges, and completion vectors used by the
coarse replay.

The outer pair $(g_\star,g_\circ)$ has law $\gamma_n^{\otimes2}$.
The internal law $\mu^{\rm int}_n$ is the fixed product law of
$(g_1,g_2,a_1,a_2,\text{thresholds})$, so
$\PP_n=\gamma_n^{\otimes2}\otimes\mu^{\rm int}_n$. After an outer pair is
fixed, the internal law remains the same. For each fixed angle,
$g_\theta$ is standard Gaussian, and $a_\theta\sim\mu^{\rm aux}_n$
independently of all matrix drivers. For the complete paths, we note that every finite collection of path
evaluations is jointly Gaussian, that the covariances within and
between blocks are unchanged by the rotation, and that cylinders at
rational times generate the path sigma algebra. After fixing $(g_\star,g_\circ)$, the spin $X_0^F$ depends only on
$(g_1,a_1,U_0^F)$, and $X_1^F$ depends only on $(g_2,a_2,U_1^F)$. We have
\begin{equation}\label{u:eq:endpoint-product-law}
 \Law_{\mu^{\rm int}_n}(X_0^F,X_1^F)
        =\nu^{\rm sm}_{g_{\rm base}}\otimes\nu^{\rm sm}_{g_{\rm base}}.
\end{equation}
This is the law of the endpoints under the internal measure, before any
Gram, pinning, or crossing event is imposed.

\begin{lemma}\label{u:lem:transfer}
If $\ex{q}<\min(\ex{k}_{\rm b},\ex{z}_{\rm P},\ex{D})$, then at each fixed
grid angle, eventually,
\begin{equation}\label{u:eq:boundary-transfer}
 \PP_n(E_{B,\theta},\ X_\theta^B\in\overline{\mathcal E}_n(g_\star))
                   \le e^{-\frac14 n^{1-\ex{q}}}.
\end{equation}
\end{lemma}

\begin{proof}
We condition only on $g_\star$. From \eqref{u:eq:coupled-rotation}, we have
\begin{equation}\label{u:eq:gaussian-regression}
 \begin{split}
 g_\theta&=(1-\delta)g_\star+\sqrt{\delta(2-\delta)}\,Z_\theta,\\
 Z_\theta&=\frac{\sqrt{1-\delta}\,g_\circ
                    +\cos\theta\,g_1+\sin\theta\,g_2}{\sqrt{2-\delta}},
 \end{split}
\end{equation}
where $Z_\theta\sim\gamma_n$ is independent of $g_\star$. For a fixed
spin, $H^c_Z(x)$ is Gaussian with variance $n/2$, including the completed
diagonal. It follows that
\begin{equation}\label{u:eq:regression-energy-moment}
 \E[e^{\beta H^c_{g_\theta}(x)}\mid g_\star]
 =e^{\beta(1-\delta)H^c_{g_\star}(x)+\beta^2n\delta(2-\delta)/4}.
\end{equation}
At a fixed angle, the auxiliary input and thresholds are independent of
the entire matrix tuple. The restricted coarse kernel of
\cref{u:prop:coarse} therefore still applies after adjoining
$g_\star,g_\circ,g_\theta$. The conditional Gaussian in
\eqref{u:eq:gaussian-regression} has positive covariance, so it does not
charge the common $\gamma_n$-null exceptional set of that proposition.

Recall that on the support of $\overline{\mathcal E}_n(g_\star)$, we have
$\|W(g_\star)\|_{\mathrm{op}}\le12$ and $|H^c_{g_\star}(x)|\le6n$.
Using \cref{u:prop:coarse} and \eqref{u:eq:regression-energy-moment}, then
summing over the set $\overline{\mathcal E}_n(g_\star)$ fixed by the target,
we bound the conditional probability of hitting this set by
\[
 \exp\bigl\{n^{1-\ex{k}_{\rm b}}+F^c(g_\star)-nP_\beta
                    +(6\beta+\beta^2/2)n\delta\bigr\}
             \pi_{g_\star}(\overline{\mathcal E}_n(g_\star)).
\]
Here we retained the target Gibbs normalization. By
\eqref{u:eq:target-pressure-event} and
\eqref{u:eq:annealed-boundary-mass}, its average in $g_\star$ is at most
\[
 \exp\{n^{1-\ex{k}_{\rm b}}+n^{1-\ex{z}_{\rm P}}
                   +C_\beta n^{1-\ex{D}}\}
                         e^{-\frac12 n^{1-\ex{q}}}.
\]
By the strict inequalities between the exponents, this proves
\eqref{u:eq:boundary-transfer}. The cost of the regression is of order
$n\delta$; the exceptional event \eqref{u:eq:target-pressure-tail} for the
target pressure is included only after the crossing
argument in \cref{u:sec:exclusion}.
\end{proof}

\begin{lemma}\label{u:lem:stability}
There exists a measurable cap event, with failure probability at most
$C(N_{\ex{f}}+2)e^{-cn}$, that controls all angles. On this event and all
coarse Gram events, neighboring grid points satisfy
\begin{equation}\label{u:eq:coarse-continuity}
 \normn{M_{j+1}-M_j}
 \le e^{C_\beta n^{\ex{a}}\log n}|\theta_{j+1}-\theta_j|
 \le e^{-n^{\ex{g}_{\rm rot}}/2}
\end{equation}
eventually, provided $\ex{a}<\ex{g}_{\rm rot}$.
\end{lemma}

\begin{proof}
We first bound the primitive inputs uniformly in the angle. We then
compare the two coarse recursions at neighboring grid points.

For a Wiener path of duration $T$, the reflection bound gives
$\PP(\sup|B|>u\sqrt T)\le4e^{-u^2/2}$. Then, for some fixed small parameter, the exponential moment of
$\sup|B|^2/T$ is bounded by a universal constant. By Chernoff's inequality
over the independent rows, we have
$\sum_i\sup|B_i|^2\le CTn$ except with probability $e^{-cn}$ when $C$
is sufficiently large. A bridge has supremum at most twice that of its
generating Wiener path. These bounds cover the two $D$ blocks, the
auxiliary cells, and the finitely many used bridge blocks. We also bound
every coordinate of the four primitive matrices and of the two streams of
completion vectors by $n$; this adds failure probability at most
$C(n^2+nN_{\ex{f}})e^{-n^2/2}$. All required path norms, completion norms,
and matrix operator norms are then at most $n^{C_0}$ for a fixed $C_0$;
for example, an enlarged $n^4$ cap suffices. Rotated inputs satisfy the
same cap up to a factor of two, and their differences are at most a
polynomial times $|\theta-\phi|$. This bounds the inputs for all angles.

For the recursion, we use the path norm
$(n^{-1}\sum_i\sup_t|b_i(t)|^2)^{1/2}$. On the target Gram events at the two grid points
$\|A_{\mathcal T}-I\|\le3/4$, every Gram--Schmidt residual has normalized
norm at least $1/2$. Indeed, the squared norm of each residual is a Schur complement
and is at least the minimum eigenvalue of the relevant Gram matrix.
If $u_l,u_l'$ are the earlier unit queries, the operator norm of their
projection difference is at most $2\sum_l\normn{u_l-u_l'}$.
Normalization is $4$-Lipschitz on vectors of norm at least $1/2$.
The next query difference is bounded by a fixed multiple of the target
difference and this sum of earlier differences.

Each target difference is bounded by $2s_j^{-1}$ times the scalar field
difference on the preceding path. Using Picard's construction, Gronwall's inequality, and integration by
parts for bounded-variation coefficients, we bound that difference by
$Ce^{V_{\ex{f}}}$ times the gap between the inputs to the preceding drivers. All $s_j^{-1}$ are polynomially bounded, since
$s_j^2\ge h/R_{V_{\ex{f}}}$ and \cref{u:lem:tails} applies. By the query
endpoint formula in \cref{u:lem:assembly}, the next endpoint gap is bounded
by a polynomial times the primitive matrix and completion gap plus the
sum of preceding query gaps. Appending the next cell adds at most the
old path gap, $\sqrt{h_j}$ times this endpoint gap, and the bridge gap.
Here $\sqrt{h_j}\le1$. The recursion implies that one chronological step multiplies the accumulated gap by
at most $n^{C_\beta}$ and adds at most that factor times the primitive
input gap. After $N_{\ex{b}}\le8n^{\ex{a}}$ steps, the
amplification is at most $e^{C_\beta n^{\ex{a}}\log n}$.
The final map $m$ is spatially $1$-Lipschitz. The grid spacing is at most
$e^{-n^{\ex{g}_{\rm rot}}}$, proving \eqref{u:eq:coarse-continuity}.
This comparison uses only coarse queries and the Gram events at the
two endpoints of each grid edge.
\end{proof}

\trilogyneedspace{12\baselineskip}
\begin{lemma}\label{u:lem:pinning}
If $\ex{r}<\ex{b}/2$, each of the $\ell_n+3$ pairs
\[
 \begin{gathered}
 (X,M_B)=(X_j^B,M_j)\quad(0\le j\le\ell_n),\\
 (X_0^F,M_0),\qquad (X_1^F,M_{\ell_n})
 \end{gathered}
\]
satisfies
\begin{equation}\label{u:eq:pinning-tail}
 \PP_n(\|X-M_B\|_2^2>R_n/4)\le2e^{-R_n/64}
\end{equation}
eventually.
\end{lemma}

\begin{proof}
Under the unconditioned law at a single angle, the source identification
in \cref{u:sec:proposal} gives the law of the complete coordinate rows,
which are independent under this law.
For $T=V_{\ex{b}}$ or $V_{\ex{f}}$, posterior projection at $S=V_{\ex{b}}$ gives
\[
 \E[(X_T-m_S)^2]=1-\E m_S^2=p(S)\le C_\beta n^{-\ex{b}/2}.
\]
The squared distances lie in $[0,4]$ and are independent across rows.
Using $e^{z/4}\le1+(e-1)z/4$ for $0\le z\le4$ and Chernoff's inequality
with parameter $1/4$, we find that, once $np(S)\le R_n/16$, their sum
exceeds $R_n/4$ with probability at most $e^{-R_n/32}$. This implies \eqref{u:eq:pinning-tail}.
For the fine endpoints, the equality of the common initial segment after
\cref{u:lem:quad} gives the same $m_S$. We have proved the bounds under the unconditioned law of each run.
\end{proof}

Define the sequence
\begin{equation}\label{u:eq:connecting-sequence}
 f_0=X_0^F,\qquad f_{j+1}=X_j^B\quad(0\le j\le\ell_n),
 \qquad f_{\ell_n+2}=X_1^F.
\end{equation}
Every step has a coarse endpoint. Let $\mathrm{Good}_n$ be the intersection
of the cap event, all $E_{B,\theta_j}$, all pinning events in
\eqref{u:eq:pinning-tail}, and the avoidance events
$X_j^B\notin\overline{\mathcal E}_n(g_\star)$ for every coarse vertex.
This is a finite intersection of measurable events on the product space.

\begin{proposition}\label{u:prop:sequence}
Assume the hypotheses of
\cref{u:prop:coarse,u:lem:boundary,u:lem:transfer,u:lem:stability,u:lem:pinning}.
On $\mathrm{Good}_n$, every step of \eqref{u:eq:connecting-sequence} has
Hamming length at most $R_n$. Moreover,
\begin{equation}\label{u:eq:sequence-union-bound}
 \begin{split}
 \PP_n(\mathrm{Good}_n^c)
 &\le C(N_{\ex{f}}+2)e^{-cn}
       +(\ell_n+1)e^{-n^{1-\ex{s}_{\rm b}}}\\
 &\quad+2(\ell_n+3)e^{-R_n/64}
       +(\ell_n+1)e^{-\frac14 n^{1-\ex{q}}}.
 \end{split}
\end{equation}
If, in addition,
\begin{equation}\label{u:eq:sequence-conditions}
 \begin{gathered}
 \ex{r}<\ex{s}_{\rm b},\qquad \ex{q}<\ex{s}_{\rm b},\qquad
 \ex{a}<\ex{g}_{\rm rot}<1-\ex{s}_{\rm b},\\
 \ex{s}_{\rm b}<\ex{s}_{\rm seq}<\ex{e},
 \end{gathered}
\end{equation}
then, eventually,
\begin{equation}\label{u:eq:sequence-failure}
 \PP_n(\mathrm{Good}_n^c)\le e^{-n^{1-\ex{s}_{\rm seq}}}.
\end{equation}
\end{proposition}

\begin{proof}
For spins $x,y$ and real vectors $m,m'$, we have
\[
 4d_H(x,y)\le3\{\|x-m\|_2^2+\|m-m'\|_2^2+\|m'-y\|_2^2\}.
\]
Each term is at most $R_n/4$ on $\mathrm{Good}_n$, eventually, by
\eqref{u:eq:coarse-continuity} and \eqref{u:eq:pinning-tail}. For the steps between fine and coarse spins, we have that the middle
term is zero, since the common initial segments agree.

A union bound using \cref{u:prop:coarse},
\eqref{u:eq:boundary-transfer}, the cap bound, and
\eqref{u:eq:pinning-tail} proves \eqref{u:eq:sequence-union-bound}.
The strict gaps in \eqref{u:eq:sequence-conditions} absorb its constants
and $\log(\ell_n+1)=O(n^{\ex{g}_{\rm rot}})$, giving
\eqref{u:eq:sequence-failure}. The fine endpoints contribute their pinning
failures to this union. The Gram and likelihood events concern only coarse proposals. Each
bound holds under the unconditioned law of the corresponding run. This union bound does not use independence between different angles.
\end{proof}

\section{Exclusion of low-flow witnesses}\label{u:sec:exclusion}

We now combine coverage with the connecting sequence. The crossing
argument is the one in \cite[Section~6, Proposition~7]{oai-typical-upper}.

\begin{proposition}\label{u:prop:witness}
Assume the hypotheses of
\cref{u:lem:smoothing,u:lem:coverage,u:prop:sequence}, and let
$0<\ex{z}_{\rm P}<1/2$. Then
\begin{equation}\label{u:eq:witness-decay}
 \begin{split}
 \gamma_n(\mathcal W_n(\epsilon_n,\phi_n))
 &\le C_\beta\bigl(n^{\ex{u}_{\rm E}-\ex{D}}
                 +n^{\ex{u}_{\rm E}-\ex{k}_{\rm f}+\ex{D}}\bigr)
        +e^{-c_\beta n^{1-2\ex{z}_{\rm P}}}\\
 &\quad+e^{2n^{1-\ex{e}}-n^{1-\ex{s}_{\rm seq}}}=o(1).
 \end{split}
\end{equation}
\end{proposition}

\begin{proof}
Fix an outer pair $(g_\star,g_\circ)\in\mathcal C_n$ with
$g_\star\in\mathcal H_n\cap\mathcal W_n(\epsilon_n,\phi_n)$.
The selected cut $S_{g_\star}$ is $\epsilon_n$-balanced. By
\eqref{u:eq:endpoint-product-law} and \cref{u:lem:coverage}, the event
$f_0\in S_{g_\star}$, $f_{\ell_n+2}\notin S_{g_\star}$ has internal
probability at least $\varpi_n$.

If a sequence joins opposite sides with steps of Hamming length at most
$R_n$, then some step has endpoints on opposite sides.
A shortest path in the hypercube between the endpoints of that step
contains an edge across the cut, and both endpoints of this edge lie in the
boundary.
Both endpoints of the step are within $R_n$ of that boundary. One endpoint
is coarse by \eqref{u:eq:connecting-sequence}, so it belongs to
$\partial_{R_n}S_{g_\star}$, which is contained in
$\mathcal E_n(g_\star)=\overline{\mathcal E}_n(g_\star)$ on this target.
Crossing is therefore impossible on $\mathrm{Good}_n$, and
\begin{equation}\label{u:eq:conditional-crossing-obstruction}
 \varpi_n\le
 \mu^{\rm int}_n\bigl(\mathrm{Good}_n(g_\star,g_\circ)^c\bigr).
\end{equation}
The crossing event may depend on the fixed outer pair, while \eqref{u:eq:endpoint-product-law} gives the law of the endpoints
under the unconditioned internal measure.

Integrating \eqref{u:eq:conditional-crossing-obstruction} against
$\gamma_n^{\otimes2}$ and using Fubini's theorem for the product measure, we obtain
\[
 \begin{split}
 \gamma_n(\mathcal W_n(\epsilon_n,\phi_n))
 &\le C_\beta\bigl(n^{\ex{u}_{\rm E}-\ex{D}}
                +n^{\ex{u}_{\rm E}-\ex{k}_{\rm f}+\ex{D}}\bigr)
        +e^{-c_\beta n^{1-2\ex{z}_{\rm P}}}\\
 &\quad+\varpi_n^{-1}\PP_n(\mathrm{Good}_n^c).
 \end{split}
\]
Here we used \cref{u:lem:smoothing} and
\eqref{u:eq:target-pressure-tail} for the outer exceptions. By
\eqref{u:eq:crossing-probability} and \eqref{u:eq:sequence-failure}, this
proves \eqref{u:eq:witness-decay}. The exceptional events for the pressure and for coverage each enter
once, before we divide by the crossing probability.
\end{proof}
